% Figure: data collapse under nondimensionalisation. Sphere drag data,
% scattered when plotted dimensionally, collapsing onto one curve when
% plotted as drag coefficient versus Reynolds number.
\begin{figure}[htbp]
\centering
\begin{tikzpicture}[font=\small]
% ===== Left panel: dimensional scatter =====
\begin{scope}[xshift=0cm]
\draw[->, thick, draw=engblack!70] (0,0) -- (5.2, 0)
  node[right, font=\scriptsize] {speed $U$ (m/s)};
\draw[->, thick, draw=engblack!70] (0,0) -- (0, 4.3)
  node[above, font=\scriptsize] {drag force $F$ (N)};
\node[font=\small, anchor=south west] at (0.2, 4.1) {\textbf{(a) dimensional}};
% Three families of points (three different sphere diameters / fluids),
% each a different quadratic-ish curve: F = c U^2 with different c.
\foreach \c/\col in {0.18/engred!85, 0.34/engblue!85, 0.09/engblack!60} {
  \foreach \U in {0.6,1.2,1.8,2.4,3.0,3.6} {
    \pgfmathsetmacro\Fy{\c*\U*\U}
    \filldraw[\col] (\U*1.3, {\Fy}) circle (0.05);
  }
}
\node[anchor=west, font=\scriptsize] at (1.4, 3.7) {\textcolor{engred!85}{$\bullet$} $d_1$};
\node[anchor=west, font=\scriptsize] at (2.4, 3.7) {\textcolor{engblue!85}{$\bullet$} $d_2$};
\node[anchor=west, font=\scriptsize] at (3.4, 3.7) {\textcolor{engblack!60}{$\bullet$} $d_3$};
\end{scope}
% ===== Right panel: nondimensional collapse =====
\begin{scope}[xshift=7.2cm]
\draw[->, thick, draw=engblack!70] (0,0) -- (5.2, 0)
  node[right, font=\scriptsize] {$\log_{10}\mathrm{Re}$};
\draw[->, thick, draw=engblack!70] (0,0) -- (0, 4.3)
  node[above, font=\scriptsize] {$\log_{10} C_D$};
\node[font=\small, anchor=south west] at (0.2, 4.1) {\textbf{(b) collapsed}};
% Ticks Re
\foreach \xt/\xv in {0.5/0, 1.6/1, 2.7/2, 3.8/3, 4.9/4} {
  \draw[draw=engblack!70] (\xt,0) -- (\xt,-0.07);
  \node[font=\scriptsize, anchor=north] at (\xt,-0.09) {\xv};
}
% Standard C_D(Re) curve (Stokes 24/Re at low Re, plateau ~0.4 in
% intermediate, drag crisis dip near Re~3e5). Plot schematically.
\draw[very thick, draw=engblue!80, smooth, samples=80, domain=0.3:4.6]
  plot (\x*1.1+0.1, { 2.5 - 1.0*\x + 0.16*\x*\x });
% All three families' points now on one curve
\foreach \dx in {0.6,1.2,1.8,2.4,3.0,3.6,4.2} {
  \pgfmathsetmacro\yy{2.5 - 1.0*\dx + 0.16*\dx*\dx}
  \filldraw[engred!85] (\dx*1.1+0.1, \yy) circle (0.045);
}
\node[anchor=west, font=\scriptsize, color=engblack!70] at (1.4, 3.5)
  {one curve: $C_D = C_D(\mathrm{Re})$};
\end{scope}
\end{tikzpicture}
\caption[Data collapse under the Buckingham Pi reduction]{Sphere-drag
measurements at three diameters and fluid combinations. (a) Plotted in
dimensional variables, the force-versus-speed data form three separate
families, one per sphere. (b) Plotted as the dimensionless drag
coefficient $C_D = F / (\tfrac{1}{2}\rho U^2 A)$ against the Reynolds
number $\mathrm{Re} = \rho U d / \mu$, the three families collapse onto a
single curve. The Buckingham Pi reduction (four dimensional variables,
three primary dimensions, hence one dimensionless group beyond $C_D$)
predicted the collapse before any datum was taken: the dimensionless drag
coefficient can depend only on the dimensionless Reynolds number. The
collapse both compresses the data and exposes the regimes (Stokes
$1/\mathrm{Re}$ at low $\mathrm{Re}$, the intermediate plateau, the
drag-crisis dip near $\mathrm{Re}\approx 3\times 10^5$).}
\label{fig:vol02:ch18:buckingham-collapse}
\end{figure}
